\chapter{1929:Arthur Harden and Hans von Euler-Chelpin}

\label{chap:chemistry1929}

The Nobel Prize in Chemistry for the year 1929 was awarded jointly to Arthur Harden and Hans von Euler-Chelpin.

\textbf{Motivation:}

for their investigations on the fermentation of sugar and fermentative enzymes

\section{Arthur Harden}

\subsection*{Early Life and Education}

Arthur Harden was born on 1865-10-12 in Manchester, England.
He was a British biochemist who shared the 1929 Nobel Prize in Chemistry with Hans von Euler-Chelpin for their investigations on the fermentation of sugar and fermentative enzymes.

\subsection*{Career and Major Research}

Harden studied chemistry at Owens College in Manchester and received his doctorate in 1894 at the University of Erlangen in Germany.
In 1897 he became head of the chemistry department at the British Institute of Preventive Medicine in London, later known as the Lister Institute, and from 1907 he also held a chair of biochemistry at the University of London.
In a long series of experiments on the fermentation of sugars by yeast-juice, Harden and his assistant William John Young showed that the addition of phosphate is essential and that the phosphate becomes bound during fermentation.
They isolated the hexose phosphate ester that carries the phosphate, later known as the Harden-Young ester (fructose 1,6-bisphosphate), and showed that the fermenting system requires a heat-stable, dialysable cofactor, which they named cozymase, in addition to the heat-labile enzyme itself.
He was knighted in 1926.

\subsection*{Nobel Prize-winning Work}

The work for which Arthur Harden was awarded the Nobel Prize in Chemistry in 1929 was described by the Nobel Committee as: "for their investigations on the fermentation of sugar and fermentative enzymes".

Harden's work clarified the role of enzymes and coenzymes in alcoholic fermentation.
He showed that fermentation proceeds through a series of phosphate-carrying intermediates and that the fermenting system consists of a heat-labile enzyme together with a heat-stable coenzyme, results that became the basis for understanding how sugar is broken down in living cells.

\subsection*{Significance and Impact}

Harden's discovery of phosphate esterification in fermentation identified the first intermediate stages of sugar metabolism and revealed the existence of coenzymes, concepts that underlie modern biochemistry.

\subsection*{Later Career and Honors}

Harden continued to work at the Lister Institute until his retirement.
He was elected a Fellow of the Royal Society in 1909 and held honorary degrees from several universities.
He died on 1940-06-17 in Bourne End, England.

\subsection*{Legacy}

The legacy of Arthur Harden endures through the lasting impact of his scientific contributions.
The Harden-Young ester and the concept of the coenzyme remain part of the basic vocabulary of biochemistry.

\subsection*{Modern Developments}

Harden's studies of yeast fermentation revealed coenzymes, the small molecules that assist enzymes, founding coenzyme research and the modern understanding of metabolism. Coenzyme chemistry underpins the entire field of bioenergetics, vitamin biochemistry, and the industrial production of biofuels and fermentation products such as ethanol and amino acids.

\subsection*{Selected Publications}

\begin{itemize}

\item Harden, A., \& Young, W. J. (1906). "The Alcoholic Ferment of Yeast-Juice." Proceedings of the Royal Society of London, Series B.

\item Harden, A. (1911). "Alcoholic Fermentation." London: Longmans, Green.

\end{itemize}

\clearpage

\section{Hans von Euler-Chelpin}

\subsection*{Early Life and Education}

Hans Karl August Simon von Euler-Chelpin was born on 1873-02-15 in Augsburg, Germany.
He was a German-born Swedish biochemist who shared the 1929 Nobel Prize in Chemistry with Arthur Harden for their investigations on the fermentation of sugar and fermentative enzymes.

\subsection*{Career and Major Research}

von Euler-Chelpin studied chemistry in Munich and Berlin and received his doctorate at the University of Berlin in 1895.
After working with Walther Nernst, he became an assistant to Svante Arrhenius at the Royal Institute of Technology in Stockholm in 1897, and he took Swedish citizenship in 1902.
From 1906 until 1941 he was professor of general and organic chemistry at the University of Stockholm, and in 1929 he became director of its new institute of biochemistry.
He studied enzymes and coenzymes with the methods of physical chemistry; in particular he investigated cozymase, the coenzyme required by zymase in alcoholic fermentation, and his laboratory contributed substantially to establishing its structure, later identified as nicotinamide adenine dinucleotide.
He also studied the vitamins together with Paul Karrer and others.
His son Ulf von Euler received the 1970 Nobel Prize in Physiology or Medicine.

\subsection*{Nobel Prize-winning Work}

The work for which Hans von Euler-Chelpin was awarded the Nobel Prize in Chemistry in 1929 was described by the Nobel Committee as: "for their investigations on the fermentation of sugar and fermentative enzymes".

von Euler-Chelpin's investigations applied the methods of physical chemistry to alcoholic fermentation and to the enzymes that bring it about.
He showed that the fermenting system consists of a thermostable coenzyme, cozymase, together with a thermolabile enzyme, and he worked out the way in which they act during the fermentation of sugar.

\subsection*{Significance and Impact}

His work established the general principle that many enzymes require a coenzyme in order to act, a concept of central importance for biochemistry.
It also connected the chemistry of fermentation to the processes by which muscles obtain energy.

\subsection*{Later Career and Honors}

von Euler-Chelpin remained active in research after his retirement and published books on enzyme chemistry, including his monograph "Chemie der Enzyme".
He held honorary doctorates from several universities.
He died on 1964-11-06 in Stockholm, Sweden.

\subsection*{Legacy}

The legacy of Hans von Euler-Chelpin endures through the lasting impact of his scientific contributions.
His clarification of the role of cozymase in fermentation is regarded as one of the foundations of modern biochemistry.

\subsection*{Modern Developments}

von Euler-Chelpin's work on fermentation enzymes and coenzymes established the enzyme chemistry central to modern biochemistry and biotechnology. His identification of cozymase (NAD) opened the field of redox biochemistry that now explains metabolism, cellular signaling, and aging research. Enzymology inspired by his work powers modern biocatalysis and biosensors.

\subsection*{Selected Publications}

\begin{itemize}

\item von Euler-Chelpin, H. (1910). "Allgemeine Chemie der Enzyme."

\item von Euler-Chelpin, H. (1915). "Chemie der Hefe und der alkoholischen Gärung." Leipzig.

\item von Euler-Chelpin, H. (1925--1934). "Chemie der Enzyme."

\end{itemize}

\clearpage

\section*{Chapter Summary}

The 1929 prize recognized Arthur Harden and Hans von Euler-Chelpin for their investigations of the fermentation of sugar and fermentative enzymes. Their discovery of the heat-stable coenzyme of alcoholic fermentation was a foundation of modern enzymology and biochemistry.
